\subsection{零半径奇层形式芽与外部 \texorpdfstring{$\pi$}{pi} 机制的局部排除}\label{subsec:conclusion-foldgauge-pi-zero-radius-external-obstruction}
沿用定理 \ref{thm:fold-bin-gauge-constant-stirling-bernoulli-hierarchy}、定理 \ref{thm:conclusion-foldgauge-pi-oddlayer-annihilation-hierarchy} 与定理 \ref{thm:conclusion-foldgauge-pi-koenigs-vs-agm-nonconjugacy} 的记号。前述链条已经把折叠 \(\pi\) 通道的奇次 Bernoulli 模态、有限窗加速器与 Koenigs 正规形完全显式化。另一方面，Borwein 型 theta/eta 模函数迭代、Ramanujan Machine 的多项式连分式以及 BBP/Bellard 的远位尾核分别代表公开文献中三类彼此不同的 \(\pi\) 机制 \cite{BorweinGarvan1997EtaPi,Wang2026RamanujanMachinePi,Bailey2023BBPCompendium}。把这些外部原型与本文的奇层谱链对照后，还可进一步抽出如下五条排除性后果：折叠 \(\pi\) 误差并不来自任何正半径解析 germ，而是一条真正的零半径 Bernoulli 奇层形式芽；因此 Borwein 型模函数加速与 Ramanujan Machine 型超几何连分式都不可能成为其局部解析正规形；同样，BBP/Bellard 与 streaming spigot 的尾核属于完全不同的代数--指数类，而 Buffon 类随机反演若要追平 fold 主阶，则必须支付二次指数样本代价。

\begin{theorem}[fold-\texorpdfstring{$\pi$}{pi} 奇层形式芽的零收敛半径]\label{thm:conclusion-foldgauge-pi-zero-radius-odd-germ}
\leanverified{paper\_conclusion\_foldgauge\_pi\_zero\_radius\_odd\_germ}
记
\[
q:=\frac{\varphi}{2},
\qquad
A_r:=\frac{B_{2r}}{r(2r-1)}\bigl(\varphi^{-1}+\varphi^{2r-3}\bigr)
\qquad (r\ge 1),
\]
并定义形式奇级数
\[
\widehat E(s):=\sum_{r\ge 1}A_r s^{2r-1}.
\]
则 \(\widehat E\) 的收敛半径为 \(0\)。特别地，
\[
\log C_m-\log(2\pi)\sim \widehat E(q^m)
\]
所定义的 fold-\(\pi\) 误差机制并不是某个解析的 \(q\)-级数，而是一条真正的零半径 Bernoulli 奇层形式芽。
\end{theorem}
\begin{proof}
由定理 \ref{thm:fold-bin-gauge-constant-stirling-bernoulli-hierarchy}，对任意固定整数 \(R\ge 1\) 都有
\[
\log C_m-\log(2\pi)
=
\sum_{r=1}^{R}A_r q^{(2r-1)m}
+O\!\bigl(q^{(2R+1)m}\bigr).
\]
因此只需研究系数 \(A_r\) 的增长。由偶 Bernoulli 数的标准渐近，
\[
|B_{2r}|
\sim
\frac{2(2r)!}{(2\pi)^{2r}},
\]
而
\[
\varphi^{-1}+\varphi^{2r-3}\asymp \varphi^{2r-3},
\]
故存在常数 \(c>0\) 与 \(r_0\ge 1\)，使得当 \(r\ge r_0\) 时
\[
|A_r|
\ge
c\,
\frac{(2r)!}{r^2}
\left(\frac{\varphi^2}{4\pi^2}\right)^r.
\]
再由 Stirling 公式，
\[
\bigl((2r)!\bigr)^{1/(2r-1)}
\sim
\frac{2r}{e},
\]
从而
\[
|A_r|^{1/(2r-1)}\ge c_1 r
\]
对充分大的 \(r\) 成立，其中 \(c_1>0\) 为常数。于是
\[
\limsup_{r\to\infty}|A_r|^{1/(2r-1)}=+\infty.
\]
而 \(\widehat E(s)\) 的非零系数恰位于奇次幂 \(s^{2r-1}\) 上，故 Cauchy--Hadamard 公式给出其收敛半径
\[
R(\widehat E)^{-1}
=
\limsup_{r\to\infty}|A_r|^{1/(2r-1)}
=+\infty.
\]
因此 \(R(\widehat E)=0\)。证毕。
\end{proof}

\begin{theorem}[Borwein--模函数加速族与折叠标量通道的局部解析不可等价性]\label{thm:conclusion-foldgauge-pi-borwein-local-analytic-obstruction}
\leanverified{paper\_conclusion\_foldgauge\_pi\_borwein\_local\_analytic\_obstruction}
设 \(F(t)\) 是任一在 \(t=0\) 解析且满足 \(F(0)=0\) 的 Borwein 型局部误差核。则不存在两个在原点收敛的幂级数
\[
\psi(s)=u_1 s+O(s^2),
\qquad
h(t)=v_1 t+O(t^2),
\qquad
u_1v_1\neq 0,
\]
使得
\[
h\!\bigl(\widehat E(s)\bigr)=F(\psi(s))
\]
作为形式级数恒等成立。换言之，Borwein 型模函数加速不可能成为折叠标量 \(\pi\) 通道的局部解析正规形。
\end{theorem}
\begin{proof}
反设存在所述 \(\psi\) 与 \(h\)。由于 \(F\) 与 \(\psi\) 都在原点解析，且 \(\psi'(0)=u_1\neq 0\)，故复合级数 \(F(\psi(s))\) 在 \(s=0\) 具有正的收敛半径。于是左端
\[
h\!\bigl(\widehat E(s)\bigr)
\]
也必须代表一个正半径解析函数。

另一方面，\(h'(0)=v_1\neq 0\)，由逆函数定理可知 \(h\) 在原点附近存在解析逆 \(h^{-1}\)。因此
\[
\widehat E(s)=h^{-1}\!\bigl(F(\psi(s))\bigr)
\]
也应在 \(s=0\) 具有正的收敛半径。这与定理 \ref{thm:conclusion-foldgauge-pi-zero-radius-odd-germ} 矛盾。故所述局部解析等价不可能成立。证毕。
\end{proof}

\begin{theorem}[Ramanujan Machine 多项式连分式类的超几何正则障碍]\label{thm:conclusion-foldgauge-pi-ramanujan-machine-hypergeometric-obstruction}
设 \(G(t)\) 是一类非典范 \(\pi/4\) 多项式连分式所对应的 contiguous \({}_2F_1\) 比值误差核，并满足 \(G(0)=0\)。则不存在两个在原点收敛的幂级数
\[
\psi(s)=u_1 s+O(s^2),
\qquad
h(t)=v_1 t+O(t^2),
\qquad
u_1v_1\neq 0,
\]
使得
\[
h\!\bigl(\widehat E(s)\bigr)=G(\psi(s))
\]
作为形式级数恒等成立。因而 Ramanujan Machine 所覆盖的这类多项式连分式簇，同样不可能提供 fold-\(\pi\) 通道的局部解析模型。
\end{theorem}
\begin{proof}
由 \cite{Wang2026RamanujanMachinePi}，这类多项式连分式可归约为 contiguous Gauss 超几何函数 \({}_2F_1\) 的比值。Gauss 超几何方程在 \(t=0\) 为常规点，因此其局部解在 \(t=0\) 邻域内解析；只要所取分母在 \(t=0\) 不为零，相应比值 \(G(t)\) 亦在原点具有正的收敛半径。

若存在所述 \(\psi\) 与 \(h\)，则与定理 \ref{thm:conclusion-foldgauge-pi-borwein-local-analytic-obstruction} 完全相同的逆函数论证会推出 \(\widehat E(s)\) 也应在 \(s=0\) 具有正收敛半径。这与定理 \ref{thm:conclusion-foldgauge-pi-zero-radius-odd-germ} 矛盾。故命题成立。证毕。
\end{proof}

\begin{theorem}[BBP/Bellard/spigot 尾核的代数--指数型与 fold 奇层谱的不相容性]\label{thm:conclusion-foldgauge-pi-bbp-tail-algebraic-exponential-obstruction}
设 \(b>1\) 为整数，\(R(x)\in\QQ(x)\) 在充分大整数上无极点，且满足 \(R(x)=O(x^{-1})\) 当 \(x\to+\infty\)。定义尾核
\[
T_m:=\sum_{k\ge m} b^{-k}R(k).
\]
则对每个整数 \(N\ge 1\)，存在常数 \(\gamma_1,\dots,\gamma_N\) 使得
\[
T_m
=
b^{-m}\left(
\sum_{\nu=1}^{N}\frac{\gamma_\nu}{m^\nu}
+O(m^{-N-1})
\right)
\qquad (m\to\infty).
\]
进一步，对任意固定窗线性组合
\[
\widetilde T_m:=\sum_{j=0}^{J}\alpha_j T_{m+j},
\]
若 \(\widetilde T_m\not\equiv 0\)，则存在最小整数 \(\nu_0\ge 1\) 与常数 \(\delta_{\nu_0}\neq 0\)，使得
\[
\widetilde T_m
=
b^{-m}\left(
\frac{\delta_{\nu_0}}{m^{\nu_0}}
+O(m^{-\nu_0-1})
\right).
\]
因此，此类尾核始终属于代数--指数型；它不可能与
\[
\log C_m-\log(2\pi)\sim\sum_{r\ge 1}A_r q^{(2r-1)m}
\]
这一纯奇层指数谱在全阶渐近类上同构。特别地，标准 BBP/Bellard 型尾核落在这一类中，而由 Euler 变换的 Leibniz 源项生成的 streaming spigot 尾核也满足同型展开。
\end{theorem}
\begin{proof}
写 \(k=m+\ell\)，则
\[
T_m
=
b^{-m}\sum_{\ell\ge 0}b^{-\ell}R(m+\ell).
\]
由于 \(R\) 是在无穷远处衰减为 \(O(x^{-1})\) 的有理函数，对任意固定 \(N\ge 1\) 都存在常数 \(c_1,\dots,c_N\) 使得
\[
R(x)=\sum_{\nu=1}^{N}c_\nu x^{-\nu}+O(x^{-N-1})
\qquad (x\to+\infty).
\]
再对每个 \(\nu\) 把 \((m+\ell)^{-\nu}\) 按 \(m^{-1}\) 作二项展开，可得存在关于 \(\ell\) 的多项式 \(P_{\nu,N}(\ell)\) 使
\[
R(m+\ell)
=
\sum_{\nu=1}^{N}m^{-\nu}P_{\nu,N}(\ell)
+O\!\bigl(m^{-N-1}(1+\ell)^{N+1}\bigr),
\]
且该余项对 \(\ell\ge 0\) 一致。代回后得到
\[
T_m
=
b^{-m}\sum_{\nu=1}^{N}m^{-\nu}
\left(\sum_{\ell\ge 0}b^{-\ell}P_{\nu,N}(\ell)\right)
+O\!\left(
b^{-m}m^{-N-1}\sum_{\ell\ge 0}b^{-\ell}(1+\ell)^{N+1}
\right).
\]
由于
\[
\sum_{\ell\ge 0}b^{-\ell}(1+\ell)^{N+1}<\infty,
\]
故可把括号内的收敛和记为 \(\gamma_\nu\)，从而得到第一组展开。

对固定窗组合，
\[
\widetilde T_m
=
\sum_{j=0}^{J}\alpha_j T_{m+j}
=
b^{-m}\sum_{j=0}^{J}\alpha_j b^{-j}
\left(
\sum_{\nu=1}^{N}\frac{\gamma_\nu}{(m+j)^\nu}
+O(m^{-N-1})
\right).
\]
再将每个 \((m+j)^{-\nu}\) 按 \(m^{-1}\) 展开，便得到
\[
\widetilde T_m
=
b^{-m}\left(
\sum_{\nu=1}^{N}\frac{\delta_\nu}{m^\nu}
+O(m^{-N-1})
\right)
\]
对任意 \(N\) 成立。若 \(\widetilde T_m\not\equiv 0\)，则取最小的 \(\nu_0\) 使 \(\delta_{\nu_0}\neq 0\)，便得到第二式。

最后，fold-\(\pi\) 误差由定理 \ref{thm:conclusion-foldgauge-pi-zero-radius-odd-germ} 给出一条纯奇层指数谱
\[
\sum_{r\ge 1}A_r q^{(2r-1)m},
\]
其中不出现任何 \(m^{-1},m^{-2},\dots\) 的代数前因子；而上式表明一切非平凡 BBP/Bellard/spigot 尾核在固定窗下都保留 \(b^{-m}m^{-\nu_0}\) 型前导项。故两类对象不可能在全阶渐近类上同构。对 Euler 变换的 Leibniz streaming spigot 尾核，同样只需把 \(R(m+\ell)\) 的有理函数展开替换为其对应的逆幂展开，证明逐字不变。证毕。
\end{proof}

\begin{theorem}[Buffon 类随机反演对 fold 主阶的二次指数代价]\label{thm:conclusion-foldgauge-pi-buffon-quadratic-exponential-cost}
\leanverified{paper\_conclusion\_foldgauge\_pi\_buffon\_quadratic\_exponential\_cost}
设某个 Buffon 类方案把 \(\pi\) 表示为
\[
p=\frac{\kappa}{\pi}\in(0,1)
\]
的倒数，其中 \(\kappa>0\) 为已知几何常数。令
\[
X_1,\dots,X_N\stackrel{\mathrm{i.i.d.}}{\sim}\mathrm{Bernoulli}(p),
\qquad
\overline X_N:=\frac1N\sum_{i=1}^{N}X_i,
\]
并定义正则化逆均值估计器
\[
\widehat\pi_N:=\frac{\kappa}{\overline X_N+N^{-1}}.
\]
则有
\[
\EE\bigl[(\widehat\pi_N-\pi)^2\bigr]=\Theta(N^{-1}).
\]
因此，若只要求均方根误差不超过 fold 主阶
\[
q^m=\left(\frac{\varphi}{2}\right)^m,
\]
则必有
\[
N\gtrsim q^{-2m}
=
\left(\frac{2}{\varphi}\right)^{2m}.
\]
换言之，Buffon 类随机反演在指数尺度上要为 fold-\(\pi\) 的主阶精度支付二次指数代价。
\end{theorem}
\begin{proof}
记
\[
f_N(x):=\frac{\kappa}{x+N^{-1}}.
\]
则 \(\widehat\pi_N=f_N(\overline X_N)\)，并且在 \(x=p\) 邻域内 \(f_N\) 光滑。由于 \(p=\kappa/\pi\)，有
\[
f_N(p)-\pi
=
\frac{\kappa}{p+N^{-1}}-\frac{\kappa}{p}
=
O(N^{-1}).
\]
同时
\[
f_N'(p)=-\frac{\kappa}{(p+N^{-1})^2}
=
-\frac{\pi^2}{\kappa}+O(N^{-1}),
\]
故 \(f_N'(p)\) 最终远离 \(0\)。对 \(f_N\) 在 \(p\) 处作 Taylor 展开，
\[
\widehat\pi_N-\pi
=
f_N'(p)\bigl(\overline X_N-p\bigr)+O\!\bigl((\overline X_N-p)^2\bigr)+O(N^{-1}).
\]
而
\[
\EE(\overline X_N-p)=0,
\qquad
\Var(\overline X_N)=\frac{p(1-p)}{N},
\qquad
\EE\bigl[(\overline X_N-p)^4\bigr]=O(N^{-2}).
\]
因此由 delta 方法，
\[
\EE\bigl[(\widehat\pi_N-\pi)^2\bigr]
=
f_N'(p)^2\Var(\overline X_N)+O(N^{-3/2})
=
\Theta(N^{-1}).
\]
于是均方根误差为 \(\Theta(N^{-1/2})\)。若要求其不超过 \(q^m\)，便必须有
\[
N^{-1/2}\lesssim q^m,
\]
等价于
\[
N\gtrsim q^{-2m}
=
\left(\frac{2}{\varphi}\right)^{2m}.
\]
证毕。
\end{proof}

\begin{conjecture}[faithful projective \texorpdfstring{$\pi$}{pi} 通道的三寄存器下界]\label{conj:conclusion-foldgauge-pi-three-register-lower-bound}
任一若要 faithful 地承载本文的 projective fold-\(\pi\) 通道，至少应同时外置三类彼此正交的寄存器：
\[
\mathfrak R_{\mathrm{bal}},
\qquad
\mathfrak R_{\mathrm{odd}},
\qquad
\mathfrak R_{\mathrm{aux}}.
\]
其中：
\begin{enumerate}
  \item \(\mathfrak R_{\mathrm{bal}}\) 负责定理 \ref{thm:conclusion-balanced-externalization-schur-optimality}、定理 \ref{thm:conclusion-binfold-critical-budget-renyi-energy}、定理 \ref{thm:conclusion-negative-mellin-two-atom-uniqueness} 与定理 \ref{thm:conclusion-balanced-twoatom-completion-capacity-branch-selection} 所固定的纤维内平衡外置、临界谐波阶梯、二原子 Mellin 指纹与 completion 选枝；
  \item \(\mathfrak R_{\mathrm{odd}}\) 负责定理 \ref{thm:conclusion-foldgauge-pi-oddlayer-annihilation-hierarchy}、定理 \ref{thm:conclusion-foldgauge-pi-affine-log-accelerator-uniqueness}、定理 \ref{thm:conclusion-foldgauge-pi-accelerator-uniform-conditioning}、定理 \ref{thm:conclusion-foldgauge-pi-selfcalibrated-minimal-window} 与定理 \ref{thm:conclusion-foldgauge-even-zeta-finite-window-extraction} 所固定的奇层消元塔、自校准、统一条件数与偶值 \(\zeta\)-有限窗抽取；
  \item \(\mathfrak R_{\mathrm{aux}}\) 负责一切 genuinely non-fold 的外部辅助结构。若该结构属于解析局部型，则被定理 \ref{thm:conclusion-foldgauge-pi-borwein-local-analytic-obstruction} 与定理 \ref{thm:conclusion-foldgauge-pi-ramanujan-machine-hypergeometric-obstruction} 排斥；若属于代数--指数尾核型，则被定理 \ref{thm:conclusion-foldgauge-pi-bbp-tail-algebraic-exponential-obstruction} 排斥；若属于随机反演型，则由定理 \ref{thm:conclusion-foldgauge-pi-buffon-quadratic-exponential-cost} 强制支付不同类别的代价。
\end{enumerate}
换言之，一旦要求 faithful，则经典 \(\pi\) 机制不能被压缩为同一个标量局部模型；它们至多在值层上给出可比较近似，而在结构层上必须裂成上述三类不可互换的寄存轴。
\end{conjecture}

\noindent
因此，fold-gauge 标量 \(\pi\) 通道在结论层面留下的并不是某个正半径模函数核、超几何连分式核、远位局域尾核或平方根统计核，而是一条零收敛半径的 Bernoulli 奇层形式芽。Borwein 型模性加速、Ramanujan Machine 型连分式、BBP/spigot 型尾核与 Buffon 型随机抽样因而分别在解析正规形、渐近类别与代价标度上与该通道严格断裂。若要继续在本文框架内部追索可能的相位式 \(\pi\) 宿主，剩余候选便只能转向双权 \(\zeta\) 的非刚性谱几何。

\endinput
